\begin{abstract}
Network digital twins promise safer and more repeatable network experimentation, but practical evidence is often fragmented across architecture proposals, emulators, topology studies, and automated control demonstrations. This dissertation designs and evaluates a scenario-driven network digital twin that uses Docker containers, deterministic routing, Linux traffic control, and a common evidence pipeline for controlled experiments. The evaluation covers bandwidth, configured one-way delay, configured packet loss, direct and multi-router topologies, a 50-node/88-link Germany50 topology, validated AI-assisted scenario generation, and real-Docker closed-loop control.

Repeated experiments show that the platform applies measurable network conditions while retaining raw and structured evidence. At configured bandwidths of 50 and 100 Mbps, five valid runs produced mean throughputs of 41.500 and 78.780 Mbps, respectively; a separate historical 20-Mbps record measured 19.2 Mbps and is not treated as a repeated cohort. Five-run delay conditions of 0, 10, 30, and 50 ms one way produced mean measured RTTs of 0.105, 27.008, 80.303, and 133.686 ms. Configured one-way loss of 0, 1, 3, and 5\% produced mean end-to-end ping loss of 0, 3, 6, and 10\%. A five-run dual-router benchmark averaged 19.220 Mbps. The complete Germany50 topology was instantiated, its 4,224-entry route plan was validated only in dry-run mode, and real traffic was measured on three representative paths rather than all pairs. The official OpenAI attempt returned HTTP 429, while a separate Qwen-compatible path succeeded after validation. Across 80 valid real-Docker RL episodes, the threshold heuristic materially outperformed the evaluated Q-learning configuration.

The results support a reproducible, extensible workflow within the evaluated configurations, while exposing limits in Docker-host fidelity, large-topology traffic coverage, provider availability, evidence preservation, and learning-based control performance.
\end{abstract}
